\selectlanguage{hebrew}
\chapter*{תקציר}

אנו חוקרים זיהוי מטרות אדפטיבי ברעש רקע לא ידוע ואולי לא-גאוסי, תוך שימוש בנתונים משניים נטולי מטרה. אנו נוקטים בגישה היברידית המשלבת מודל ולמידה מנתונים: הלמידה משמשת למידול הרקע, בעוד שהגלאי והנרמול שלו נשמרים בצורה המוכתבת על ידי תורת הגילוי הקלאסית. אנו מתמקדים בגלאי החזק ביותר מקומית, התלוי ברקע רק דרך פונקציית הסקור שלו ומספר קטן של מומנטים של הסקור. אנו מנצלים מבנה זה על ידי אמידת מאפייני הרקע באמצעות דינוייזינג סקור מאטצ'ינג והצבת האומד בגלאי. הפרדה זו מאפשרת לעשות שימוש חוזר במודל רקע נלמד יחיד עבור חתימות מטרה שונות. במקרה הגאוסי האפיני, השיטה המוצעת מצטמצמת למסנן מותאם אדפטיבי עם רגולריזציה, מה שמסביר את יתרונותיה במשטר של מדגמים סופיים. בנוסף, אנו מפתחים הרחבה אדפטיבית מרחבית המשפרת את הדיוק תוך שמירה על כיול אמפירי של קצב אזעקות שווא קבוע בסביבות המשתנות במרחב. ניסויים על שלוש סצנות היפרספקטרליות אמיתיות ממחישים את יתרונות הגלאים המוצעים ביחס לשיטות ייחוס גאוסיות, תערובות גאוסיות, מבחן ראו נלמד, ושיטות למידה עמוקה עדכניות.
\selectlanguage{english}